\SourceLeaf{168}schlagung erwartet werden, und unter Deficit den Unterschied, wenn für die Ausgabe eine grössere Summe sich ergibt, als für die Einnahme. In einem solchen Zusammenhange sind Überschüsse und Deficit einander gerade entgegengesetzt, und das eine schliesst das andre von selbst aus.

Von einer solchen Rechnung unterscheidet sich diejenige, welche in der zweiten Abtheilung der Denkschrift fur die Witwenkasse geführt ist, in zwei wesentlichen Stücken.

1) Die letztere vergleicht nicht die Ausgaben und Einnahmen für Ein Jahr, oder für einige Jahre, sondern umfasst beide für die ganze unbegrenzte Zukunft, so weit und so genau, als es möglich ist, dieselbe im Voraus zu veranschlagen.

2) Sie summirt nicht die Gelder selbst, in der Grösse wie sie werden verausgabt oder vereinnahmt werden (was auch wegen der Unbegrenztheit ohne Sinn sein würde), sondern deren nach bestimmtem Zinsfuss auf den Anfangszeitpunkt discontirte oder reducirte Werthe. In diesem Sinne lässt sich auch eine ohne Begrenzung fortlaufende Reihe von Ausgaben oder Einnahmen doch zu einem bestimmten endlichen Resultate summiren, und diese Möglichkeit leuchtet leicht ein, wenn man bedenkt, dass eine Geldsumme durch Discontirung für einen sehr langen Zeitraum ganz enorm zusammenschmilzt, wie z. B. eine nach 422 Jahren zu leistende Zahlung von 7000 Rthl. nach dem Zinsfusse von \(3\tfrac{1}{2}\) Procent jetzt nur den Werth von \textit{einem Pfennig} hat. Noch leichter kommt man zu demselben Resultate durch die Erwägung, dass der jetzige Geldwerth einer ohne Aufhören jährlich in gleicher Grösse zu leistenden Zahlung nichts anderes ist, als die Kapitalsumme, welche nach dem gewählten Zinsfusse jährlich einen eben so grossen Zins abwirft.

Es ist nun zwar schon von selbst klar, dass wenn nach den factischen Grundlagen einer solchen Rechnung die Ausgaben in späterer Zeit grösser sein werden als jetzt. die ausser den Kapitalzinsen aber noch Statt findenden Einnahmen hingegen in der Rechnung nicht als steigend angenommen werden dürfen, das Resultat der Rechnung ein Deficit sein wird, falls jetzt Einnahme gegen Ausgabe keinen Überschuss gibt. Allein es lässt sich nicht ohne Weiteres umgekehrt behaupten, dass wenn jetzt Überschüsse von der Einnahme gegen die Ausgabe Statt finden, \textit{kein} Deficit in der Totalrechnung sein werde, denn dazu wird nicht nur das Dasein von Überschüssen sondern eine hinlängliche Grösse derselben erfordert. Hat eine richtige Total-Bilanz-Rechnung als Endresultat ein Deficit ergeben, so steht dadurch fest, dass die jetzigen jährlichen Überschüsse zu klein sind, um die in späterer Zeit bevorstehenden Ausfälle zu decken. Die richtige Rechnung hat die gegenwärtigen zeitweiligen Überschüsse schon als Bestandtheile der Bilanz mit berücksichtigt, und dieselben dürfen der Anstalt nicht \textit{zweimal} in Einnahme gestellt werden.

Eben so falsch, wie die Einbildung, dass jetzige zeitweilige Überschüsse ein jetzt vorhandenes Deficit vermindern werden, ist die Vorstellung, dass dies Deficit dann getilgt sein werde, sobald das Vermögen eine dem Betrage des Deficits gleichkommende Vergrösserung erhalten habe. Ein solcher Schluss ist nur in dem einzigen Falle zulässig, wenn die Vergrösserung des Vermögens \textit{sogleich} und zwar durch \textit{fremde} in der Bilanz nicht schon enthaltene Zuflüsse bewirkt wird. Gesetzt z. B. das Vermögen der Witwenkasse habe sich nach 20 Jahren (etwa theilweise unter Mitwirkung von neuen Zuflüssen die aber nicht fortdauernder Art wären) um 17520 Rthl. vergrössert, so wird dann alsdann doch das Deficit nicht gehoben sein, sondern es kann dann möglicherweise grösser sein als jetzt. Die Bilanz am 1. Oct. 1865 wird nemlich hervorgehen aus Vergleichung des dann vorhandenen Kassenvermögens mit dem auf \textit{diesen} Zeitpunkt discontirten Betrage aller von da an bevorstehenden Ausgaben, welcher Betrag viel grösser sein wird, als der ähnliche Betrag für den 1. Oct. 1845, und zwar aus dem einfachen Grunde, weil 1865 der Zeitpunkt der viel grössern jährlichen Ausgaben so viel näher gerückt sein wird, ja höchst wahrscheinlich diese alsdann schon in bedeutendem Maasse eingetreten sein werden.
